\documentclass[10pt]{article}

\usepackage[utf8]{inputenc}

\usepackage[T1]{fontenc}

\usepackage{amsmath}

\usepackage{amsfonts}

\usepackage{amssymb}

\usepackage[version=4]{mhchem}

\usepackage{stmaryrd}



\begin{document}

XAATA - APAKИ AKСИОМА - см. Аксиоматическая квантовая теория поля.



XAAPA MEPA на локально компактной групп е - левая (правая) мера на группе $G$, инвариантная относительно левых сдвигов $L_{h}$ :



\[

g \rightarrow h g \text { на } G

\]



(соответственно, правых сдвигов $R_{g}: g \rightarrow g h$ ). Эти две меры называются соответственно ле вой и правой ме рой Х а а p a. X. М. единственна с точностью до пропорциональности. Обычно эту меру удается задать с помощью явных формул. Напр., на группе Ли для этого достаточно задать инвариантную относительно переносов форму объема.



См. также Инвариантное интегрирование. Ю. А. Неретин. XАAPA СИСТЕМА - см. Ортогональная система функций. ХАНКЕЛЯ ФУНКЦИИ - см. Ганкеля функции.



ХAOC (в квантовой механике), квантовый хаос, - совокупность свойств, специфически присушая кван товым хаоти ческим с истемам - гамильтоновым системам, которые в классическом случае обладают стохастическим движением. Последнее означает, что в фазовом пространстве сушествуют области финитного движения, в к-рых зависимость динамич. переменных ( $(x)$ от времени $\boldsymbol{x}(t)$ экспоненциально неустойчива при изменении начальных условий:



\[

\frac{\partial x(t)}{\partial x(0)} \sim \exp (\sigma t) .

\]



Величина $\sigma$, определяющая степень такой неустойчивости, называется показателем Ляпунова (см. [1], [2]).



Квантовый X. доступен определению только в квазиклассич. случае: для систем, у к-рых характерный масштаб действия $S$ велик в сравнении с постоянной Планка $\hbar$.



Квантовый Х. имеет два основных аспекта. Первый проявляется в характеристиках собственно квантовых параметров (энергетич. спектров, волновых функций и матричных элементов операторов) и фиксирует отличия квантовых хаотич. систем от квантовых регулярных (полностью интегрируемых) систем. Второй аспект проявляется в отличии динамики наблюдаемых для квантовых хаотич. систем и ее классич. аналога. В классич. механике сушествование стохастич. движения является исключительным свойством нелинейных неинтегрируемых систем с числом степеней свободы $N \geqslant 2$. В силу линейности уравнения Шредингера истинная, асимптотическая при $t \rightarrow \infty$ экспоненциальная неустойчивость движения средних значений наблюдаемых при изменении начальных условий в квантовых системах невозможна. С другой стороны, в силу принципа соответствия значения динамич. переменных, найденные по классич. и квантовой теориям при сходных начальных условиях, должны оставаться близкими в течение нек-рого времени соответствия $\tau_{C}$, расходясь по прошествии этого времени. О различии в динамике классич. и квантовых наблюдаемых при временах $t>\tau_{C}$ говорят как о квантовом подавлении классич. стохастичности. Для данной системы время соответствия $\tau_{C}$ растет вместе с увеличением отношения $S / \hbar$.



\section{XAOC}

Проблема квантового хаоса в форме вопроса о способе квантования классич. неинтегрируемых систем известна с первых лет существования квантовой теории. Она не могла быть решена до становления теории стохастичности в классич. гамильтоновых системах, выявившей сложное, квазислучайное поведение динамич. переменных при стохастич. движении и подготовившей перенос акцента с определения конкретных значений квантовых параметров на определение их статистич. свойств (см. [1], [2]). Начало современному этапу изучения квантового Х. было положено в [3], где на основе принципа соответствия описаны качественные свойства спектров и волновых функций квантовых хаотич. систем, и в [4], где было предложено исследовать статистич. свойства энергетич. спектра и предугадано явление отталкивания уровней в квантовых хаотич. системах; в [5] теоретически определен вид статистики спектра для регулярных квантовых систем; в [6], где было обнаружено в численном эксперименте и истолковано проявление времени соответствия $\tau_{C}$ для модели квантовой хаотич. системы.



Статистические свойства параметров квантовых хаотических систем. Основные понятия, используемые для описания квантового X., заимствованы из статистич. теории случайных матриц, развитой в 60-е гг. 20 в. для описания структуры спектров сложных ядер (см. [7]). Центральное место в этой теории занимает понятие об ансамбле матриц - гамильтонианов, обладающих статистически независимыми элементами, функции распределения к-рых остаются инвариантными при преобразованиях, совместимых с дискретными симметриями гамильтониана.



Простейшей характеристикой структуры энергетич. спектра - последовательности значений уровней энергии $E_{n}-$ является функция распределения относительной величины межуровневых расстояний $S_{n}$ :



\[

S_{n}=\left(E_{n}-E_{n-1}\right) \rho\left(E_{n}\right),

\]



где $\rho(E)$ - плотность уровней; в квазиклассич. случае такая плотность хорошо определена и плавно зависит от $E$. Функция распределения $P(S)$ случайной величины $S$ нормирована условиями



\[

\int_{0}^{\infty} P(S) d S=1, \quad \int_{0}^{\infty} P(S) S d S=1

\]



и характеризует свойства спектра на минимальном масштабе энергии.



Другой важной характеристикой статистич. свойств спектра является спектральная жесткость $\Delta_{3}(L)$, к-рая определяется ниже. Для последовательности уровней $\varepsilon_{n}$, нормированных на единичную плотность $\left(\varepsilon_{n}=\varepsilon_{n-1}+S_{n}\right)$, вводится ступенчатая функция $n(\varepsilon)$, равная числу уровней с $\varepsilon_{n} \leqslant \varepsilon$. По построению $n(\varepsilon)$ имеет вид лестницы с единичным в среднем наклоном. Функция $\Delta_{3}(x, L)$ определяется ка́к минимум квадратичного отклонения $n(\varepsilon)$ на интервале $(x, x+L)$ от прямой линии:



\[

\Delta_{3}(x, L)=\frac{1}{L} \min _{A, B} \int_{x}^{x+L}(n(\varepsilon)-A \varepsilon-B)^{2} d \varepsilon .

\]



Значение $\left\langle\Delta_{3}(x, L)\right\rangle_{x}$, усредненное по значениям $x$ из области, в к-рой характер флуктуаций спектра можно считать неизменным, зависит только от $L$ и обозначается $\Delta_{3}(L)$.





\end{document}